\chapter{Conclusion and Future Work}
\label{chap:conclusion}

\section{Summary of Achievements}
\label{sec:summary}

This project set out to build a video sharing platform whose transcoding subsystem runs on a
serverless container and event-driven architecture, and to evaluate whether that choice
delivers the cost and scalability benefits it promises.

The functional platform was completed. Users can register, authenticate, upload videos directly
to Amazon S3 through pre-signed URLs, watch them with adaptive bitrate playback across three
renditions, control whether each video is public or private, manage their channel, search and
filter the catalogue, and interact through likes and comments.

The transcoding pipeline operates without manual intervention. An upload emits an S3 event that
is buffered in Amazon SQS and dispatched by an AWS Lambda function to AWS Batch, where a
container packaging FFmpeg produces the renditions, the master playlist, and a thumbnail before
terminating. The heartbeat pattern prevents duplicate processing of long jobs, and failed jobs
are retried and ultimately isolated in a dead letter queue rather than being lost silently.

Content protection was implemented in two layers. Origin Access Control keeps the storage
bucket entirely private, and CloudFront Signed Cookies ensure that private videos are rejected
at the edge even when the exact manifest URL is known. This closes a gap that would otherwise
have made the private visibility setting ineffective in practice, since the API can restrict its
own responses but cannot restrict a content delivery network.

The infrastructure is fully described as code in eleven Terraform modules across two
environments, and seven GitHub Actions workflows automate linting, unit testing, static security
analysis, container scanning, infrastructure validation, and deployment. The security gate
blocks merges on fixable critical findings rather than merely reporting them.

The automated test suite grew to 104 tests across nine suites. Its practical value was demonstrated
when it caught a defect in the cookie signing implementation that would have prevented every
private video from playing once deployed. A second, related defect in that same subsystem, caught
only later and by a different means, is discussed in Section~\ref{sec:limitations} below.

\section{Limitations}
\label{sec:limitations}

Several limitations should be stated plainly.

The bitrate ladder is fixed at three rungs regardless of content. As the literature reviewed in
Section~\ref{sec:core-models} shows, the optimal ladder is content-dependent, and a fixed ladder
therefore wastes bandwidth on simple content while under-serving complex content. Netflix
reported bitrate reductions between 28\% and 38\% at equivalent perceptual quality by optimising
the ladder per shot \cite{netflix2018dynamic}, which indicates the magnitude of the opportunity
forgone.

Each video is transcoded by a single container, without the chunk-based parallelism that
Netflix employs. Processing time consequently scales linearly with video duration, and a long
video cannot be accelerated by adding capacity. This is acceptable at the scale of this project
but would become the dominant constraint for longer content.

Signed cookies remain valid for two hours after issuance. If an owner switches a video from
public to private, viewers holding unexpired cookies retain access until those cookies lapse.
Immediate revocation would require rotating the signing key or introducing an edge function that
performs a per-request check.

A more fundamental limitation surfaced during the preparation of the experimental results for
this report. Verifying the deployed environment directly through the AWS command line interface
revealed that no CloudFront distribution had ever been created in the development AWS account,
and that the processed content bucket carried a manually added policy granting anonymous
\texttt{s3:GetObject} access to any object. In other words, the content protection scheme
described in Section~\ref{sec:content-protection} was fully implemented and validated in code
but had never actually been deployed; every video, public or private, was directly retrievable
from Amazon S3 without any authorisation check. Attempting to apply the Terraform configuration
to correct this failed with an access-denied error, because AWS requires manual account
verification before a new account may create CloudFront distributions, a restriction the project
had encountered before and for which the \texttt{enable\_cloudfront} flag had originally been
introduced as a workaround. A support request was filed with AWS to lift this restriction and, at
the time, remained pending with no bound on how long it might take. This episode is reported here
in the interest of honesty rather than omitted, and it illustrates a broader lesson about
serverless and managed cloud services: correctness of the Infrastructure as Code is necessary but
not sufficient, since provisioning can still be blocked by account-level policy outside the
application's control, and verifying claims against the live environment rather than against the
source code alone is essential before a security control can be considered in effect.

Rather than wait indefinitely on the support ticket, the gap was closed by provisioning the
distribution in a second AWS account that was not subject to the same restriction. Terraform
supports this without duplicating any module: the distribution, its Origin Access Control, its
cache and response-header policies, and its Signed Cookie trusted key group were given an explicit
\texttt{provider} argument pointing at the second account's credentials, while the buckets, Batch,
Lambda, and the database connection remained where they already were. Only one resource could not
simply move with the rest: a bucket policy can only be attached by the account that owns the
bucket, so the statement granting the distribution read access was kept on the original account's
provider. This is not a special case requiring extra logic, because an Origin Access Control
policy is already expressed as a service principal together with an \texttt{AWS:SourceArn}
condition naming the distribution, a form that is indifferent to which account owns the ARN it
names. With the cross-account distribution in place, the manually added public bucket policy was
removed, and the bucket returned to being reachable only through CloudFront, exactly as
Section~\ref{sec:content-protection} specifies. The workaround changes nothing about the security
property being evaluated: the distribution, key group, and signing logic are identical to what
would have been provisioned in the original account, and the account boundary a distribution
happens to sit behind is invisible to a viewer.

A second, more subtle defect emerged only once Signed Cookies were exercised against the newly
deployed distribution, and it is worth recording because of how it evaded the existing test suite.
The function constructing the resource pattern a cookie's policy authorises,
\texttt{buildResourcePattern} in \texttt{cloudfrontService.js}, signed
\texttt{videos/*/\{videoId\}/*}, with a wildcard segment between the fixed prefix and the video
identifier. The transcoder, however, writes every rendition and the master playlist directly
under \texttt{videos/\{videoId\}/}, with no such intermediate segment. Every cookie the backend
issued was consequently well-formed and correctly signed, and CloudFront rejected every one of
them, because the resource each cookie authorised never matched any URL a viewer could actually
request; every private and public video alike returned \texttt{403 Forbidden}. The existing test
for \texttt{buildResourcePattern} did not catch this because its assertion compared the function's
output against a literal string obtained by reading the function itself, not against the
transcoder's actual key layout, so the test could only confirm that the function had not changed,
not that the pattern it produced was one CloudFront would ever match. This is the canned-policy
defect from Section~\ref{sec:testing} recurring in a different guise: a test written by copying an
implementation's current behaviour cannot detect a defect in that behaviour, only a later
regression away from it. The fix corrects the pattern to \texttt{videos/\{videoId\}/*} and
rewrites the test to assert against the key layout the transcoder actually produces.

Two further consequences of finally exercising the deployed distribution are worth recording more
briefly. Gating every object behind Signed Cookies also gated video thumbnails, which a catalogue
listing must render for many videos at once even though a cookie authorises exactly one; the fix
carves out a separate, unsigned cache behaviour for the thumbnail path specifically, since a
thumbnail is not the content the mechanism is meant to protect. Separately, the frontend's own
CloudFront distribution, provisioned around the same time for the unrelated reason that an Amazon
S3 website endpoint cannot return \texttt{200} for a single-page application's client-side routes
on a direct navigation, had no cache-invalidation step in its deployment workflow, so a fix already
synced to the origin bucket continued to be served from the edge cache for hours afterward; the
workflow now invalidates that distribution immediately after every deployment. Both were found the
same way as the two defects above: by exercising the deployed system rather than trusting that a
correct configuration change had taken effect once it merged.

Signed Cookies were verified end to end after these corrections: a request for a manifest without
a valid cookie returns \texttt{403 Forbidden} at the edge, and the identical request succeeds once
the backend's playback-authorisation endpoint has issued cookies for that video, confirming the
mechanism evaluated in Section~\ref{sec:content-protection} is now in effect on the deployed
system and not merely present in its source code. Cloudflare continues to sit in front of the
deployment as its DNS provider and edge TLS terminator, a role introduced as a stopgap while the
distribution above did not yet exist; that role persists, but it now proxies to the real
CloudFront distribution rather than substituting for it.

The system implements no digital rights management. An authorised viewer can capture the
decrypted stream or share their cookies, and preventing this is outside the scope of the
techniques used here.

Finally, the rate limiting and view deduplication mechanisms hold their state in process memory.
This is correct for the current single-instance deployment but would need to move to a shared
store such as Redis before the backend could be scaled horizontally, since each instance would
otherwise enforce its own independent limits.

A defect found late in the evaluation deserves recording, both because it invalidated part of a
measurement and because the way it evaded detection is instructive. The Terraform module for
Secrets Manager created the entry holding the notification e-mail password unconditionally, but
populated its value only when an e-mail password had actually been configured. Leaving e-mail
unconfigured therefore produced an empty secret: a container with no version carrying the
\texttt{AWSCURRENT} label. The Batch module, meanwhile, decided whether to declare that secret in
the job definition by testing whether its ARN was non-empty --- and the ARN of an empty secret is
perfectly non-empty. Every transcoding task consequently failed at
\texttt{ResourceInitializationError} before its entrypoint ran, and the transcoding pipeline was
inoperative for as long as that job definition revision was active. The fix was to make the secret
itself conditional so that the unconfigured state is internally consistent: no secret, an empty
ARN, and no declaration in the job definition.

What makes this worth reporting is why it went unnoticed. The only jobs submitted while the defect
was live were those of the concurrency stress test, whose synthetic payloads were \emph{expected}
to fail; a wall of failed jobs was precisely what the experiment predicted, so the failures were
read as confirmation rather than as a symptom. The lesson generalises beyond this project: when an
experiment is designed such that failure is the anticipated outcome, failure carries no
information, and the reason for each failure must be inspected rather than assumed. Verifying the
\emph{cause} recorded in \texttt{statusReason}, not merely the terminal state, would have surfaced
the defect within minutes instead of hours.

\section{Future Work}
\label{sec:future-work}

Three measurements should be repeated before the experimental results are considered settled. The
Time-to-First-Frame measurement in Section~\ref{sec:qoe-evaluation} was taken against the S3 origin
before the CloudFront distribution described in Section~\ref{sec:limitations} existed, and should
be repeated against the live distribution with Signed Cookies attached exactly as a real playback
session presents them, so that the reported figures reflect what a viewer experiences on the
deployed system rather than origin latency alone. The concurrency test re-run with genuine video
payloads (Section~\ref{sec:real-throughput-validation})
confirmed that every job succeeds under load, but it was not instrumented with the same
twenty-five-second polling cadence as the original synthetic run, so a precise drain-rate figure
for real transcoding throughput --- as opposed to confirmation that throughput is non-zero and
failure-free --- remains to be measured; repeating the same short-clip payload with that polling
protocol in place would close this. The vCPU scaling comparison should also be repeated with more
replications of the 1~vCPU baseline and with memory held constant, since
Section~\ref{sec:pipeline-evaluation} showed the present estimate to be dominated by baseline
variance and confounded by a simultaneous change in memory allocation.

The most valuable functional extension would be per-title bitrate ladder optimisation driven by a
perceptual quality metric such as VMAF. Encoding a short probe of the source at several
operating points and selecting those on the convex hull of the quality-versus-bitrate curve
would reduce delivery bandwidth at unchanged perceived quality, which lowers the largest
remaining cost component of the system.

Chunk-based parallel transcoding would address the linear scaling of processing time. Splitting
the source at keyframe boundaries, encoding the segments concurrently across multiple
containers, and concatenating the results would convert additional capacity into reduced
latency. The main challenge is maintaining rate control consistency across chunk boundaries so
that quality does not fluctuate visibly.

Hardware acceleration offers a further avenue. Research on GPU-enabled serverless platforms for
HEVC encoding \cite{gpuserverless2024} suggests substantial throughput gains, and adopting a
more efficient codec such as HEVC or AV1 alongside H.264 would reduce bandwidth for clients that
support it.

Finally, moving the rate limiting and deduplication state into a shared store, and adding
distributed tracing across the S3, SQS, Lambda, and Batch boundaries, would prepare the system
for horizontal scaling and make failures in the asynchronous pipeline substantially easier to
diagnose.
